% 3-5-evaluation.tex
%  Budget: 4pp.  Converted from thesis/drafts/3-5-evaluation-framework.md,
%  with three additions the draft did not carry:
%   (a) the CONFIRMATORY vs EXPLORATORY family declaration -- mandatory,
%       HANDOFF 16A/A1. Without it, 4-5-1 cannot quote p=0.0226.
%   (b) the two-naives explanation -- HANDOFF 16A/A2 (why naive rMAE = 0.849
%       is not a contradiction).
%   (c) the sMAPE-vs-MAPE consistency argument -- HANDOFF 16A/A3, including
%       "no conclusion rests on sMAPE".
%  Source map: walk-forward protocol, 1092-day calibration window, two-year
%   held-out test, and the rMAE normalizing naive -> lago_forecasting_2021 ;
%   MAPE unreliable on markets with near-zero/negative prices, reported and
%   flagged in the reference table itself -> lago_forecasting_2021 .
%  Negative-price counts verified against data/raw/DE.csv: 297 train, 241 test.

\section{چارچوب اعتبارسنجی و معیارهای ارزیابی}
\label{sec:3-5}

\subsection*{اعتبارسنجی پیش‌رونده}

در پیش‌بینی قیمت برق، ترتیب زمانی داده‌ها خودْ اطلاعات است. هر شکلی از تقسیم تصادفی
داده به مجموعه‌های آموزش و آزمون، اطلاعات آینده را به گذشته نشت می‌دهد و خطای مدل را
به‌صورتی غیرواقع‌بینانه کم برآورد می‌کند. از این رو در این پژوهش، مطابق پروتکل مرجع
لاگو و همکاران \cite{lago_forecasting_2021}، منحصراً از اعتبارسنجی پیش‌رونده با مبدأ
غلتان استفاده شده است.

سازوکار چنین است: برای هر مبدأ پیش‌بینی، یعنی روز \lr{D}، مدل تنها بر پنجره‌ی واسنجیِ
انتهایی شامل ۱۰۹۲ روز متوالیِ منتهی به روز پیش از مبدأ آموزش می‌بیند و بردار
۲۴مؤلفه‌ای قیمت‌های ساعتی روز \lr{D+1} را پیش‌بینی می‌کند؛ سپس مبدأ یک روز جلو می‌رود
و کل فرایند تکرار می‌شود. طول پنجره‌ی واسنجی \lr{—} ۱۰۹۲ روز، معادل سه سال ۳۶۴روزه
\lr{—} مقدار پیش‌فرضِ خودِ پیاده‌سازی مرجع است و نه انتخابی از سوی نگارنده؛ همین
هم‌ارزیِ دقیق با پروتکل منتشرشده، مبنای منصفانه‌بودنِ مقایسه‌ی فصل چهارم با اعداد مرجع
است.

دوره‌ی آزمون، مطابق همان مرجع، دو سال کامل ۲۰۱۶ و ۲۰۱۷ است: ۷۲۸ مبدأ پیش‌بینی و
\lr{17{,}472} ساعت به‌ازای هر مدل.

\subsection*{جداسازی سخت‌گیرانه‌ی پنجره‌ی اعتبارسنجی از آزمون}

برای تنظیم ابرپارامترها (۵۰ آزمایش \lr{Optuna} به‌ازای هر مدل؛ \cref{sec:3-7}) پنجره‌ای
۳۶۴روزه از انتهای داده‌ی آموزش جدا می‌شود. قید بنیادین آن است که این پنجره اکیداً پیش
از آغاز دوره‌ی آزمون پایان یابد.

این قید صرفاً یک قرارداد نیست. در کد با تابعی تصریحی اعمال می‌شود که در صورت هر
تخطی \lr{—} از جمله پنجره‌ی تهی، که قابل راستی‌آزمایی نیست \lr{—} اجرای برنامه را
متوقف می‌کند. آزمون یکپارچه‌ای نیز با مقادیر واقعیِ پیکربندی می‌سنجد که حتی یک روز از
دوره‌ی آزمون در پنجره‌ی آموزشِ نخستین مبدأ حضور نداشته باشد. بدین ترتیب ادعای
«اعتبارسنجی اکیداً پیش از آزمون» در این پایان‌نامه ادعایی آزموده است، نه فرضی ضمنی.

\subsection*{سه مصنوعِ مدل، و اینکه کدام عددی را تولید می‌کند}

مخزن این پژوهش سه دسته مصنوعِ مدل نگه می‌دارد که کاربردشان کاملاً متفاوت است. چون یکی
از آن‌ها بر پنجره‌ای آموزش دیده که دوره‌ی آزمون را در بر می‌گیرد، تفکیک آن‌ها باید
پیش از هر عدد دقتی و به‌صراحت انجام شود؛ در غیر این صورت خواننده حق دارد نشت اطلاعات
را فرض بگیرد.

\textbf{یکم: مدل‌های چارچوب پیش‌رونده \lr{—} تنها منبع اعداد فصل چهارم.} در این
چارچوب، به‌ازای هر یک از ۷۲۸ مبدأ، مدل بر پنجره‌ی ۱۰۹۲روزه‌ای برازش می‌شود که اکیداً
پیش از همان مبدأ پایان می‌یابد، و تنها روز بعد از آن را پیش‌بینی می‌کند. \emph{تمام}
معیارهای دقت و \emph{تمام} آزمون‌های معناداریِ فصل چهارم از همین مسیر می‌آیند. هیچ عدد
گزارش‌شده‌ای در فصل چهارم به مدلی تکیه ندارد که روز مورد پیش‌بینی را دیده باشد.

\textbf{دوم: مدل تثبیت‌شده بر کل داده \lr{—} هیچ عددی از آن گزارش نمی‌شود.} این مصنوع
بر پنجره‌ی ۱۰۹۲روزه‌ی منتهی به ۳۱ دسامبر ۲۰۱۷ برازش شده است، یعنی پنجره‌ای که دوره‌ی
آزمون را \emph{در بر می‌گیرد}. به همین دلیل این مدل هرگز برای تولید هیچ معیار دقتی به
کار نمی‌رود. کاربرد آن تنها ابزار عملیاتیِ \cref{sec:5-3} است، که برای پیش‌بینی فردا باید از
تازه‌ترین داده‌ی موجود استفاده کند و اساساً در وضعیت ارزیابی قرار ندارد.

\textbf{سوم: مدل تفسیر \lr{—} برازشی جداگانه، اکیداً پیش از دوره‌ی آزمون.} تحلیل
\lr{SHAP} در \cref{sec:4-6} نه بر مدل‌های چارچوب پیش‌رونده انجام می‌شود و نه بر مدل تثبیت‌شده
بر کل داده. دلیل کنارگذاشتن مورد دوم روشن است: توضیح‌دادنِ روزهایی که مدل آن‌ها را در
آموزش دیده، تفسیری درون‌نمونه‌ای است. پس \cref{sec:4-6} مدلی مستقل را تفسیر می‌کند که پنجره‌ی
آموزشش هم‌طول همان ۱۰۹۲ روز است اما اکیداً پیش از مرز دوره‌ی آزمون پایان می‌یابد، و
بنابراین هر روزی که توضیح داده می‌شود واقعاً دیده‌نشده است. این تضمین در کد زندگی
می‌کند و با آزمون خودکار سنجیده می‌شود؛ پیاده‌سازی حتی از برازش بر پنجره‌ای کوتاه‌تر
سر باز می‌زند، تا این همزاد در سکوت از مدل گزارش‌شده فاصله نگیرد.

خلاصه‌ی قابل وارسی: اعداد از مصنوع نخست می‌آیند، ابزار از مصنوع دوم، و نمودارهای
تفسیرپذیری از مصنوع سوم. تنها مصنوعی که دوره‌ی آزمون را دیده است، هیچ عددی در این
پایان‌نامه تولید نمی‌کند.

\subsection*{معیارهای ارزیابی}

چهار معیار دقت گزارش می‌شود. برای پیش‌بینی $\hat{p}_i$ و مقدار واقعی $p_i$ بر
$N$ نقطه:

\begin{equation}
\mathrm{MAE} = \frac{1}{N}\sum_{i=1}^{N}\left|p_i - \hat{p}_i\right|
\end{equation}

\begin{equation}
\mathrm{RMSE} = \sqrt{\frac{1}{N}\sum_{i=1}^{N}\left(p_i - \hat{p}_i\right)^2}
\end{equation}

\begin{equation}
\mathrm{sMAPE} = \frac{100}{N}\sum_{i=1}^{N}
\frac{\left|p_i - \hat{p}_i\right|}{\left(|p_i| + |\hat{p}_i|\right)/2}
\end{equation}

\begin{equation}
\mathrm{rMAE} = \frac{\mathrm{MAE}}{\mathrm{MAE}_{\mathrm{naive}}}
\end{equation}

که در آن‌ها $N$ شمار نقاط ارزیابی، $p_i$ قیمت محقق‌شده، $\hat{p}_i$ قیمت
پیش‌بینی‌شده، و $\mathrm{MAE}_{\mathrm{naive}}$ خطای مطلق میانگینِ یک پیش‌بینیِ ساده‌ی
مرجع است. مقدار \lr{rMAE} کوچک‌تر از یک یعنی مدل از آن مرجع بهتر است و مقدار بزرگ‌تر
از یک یعنی بدتر.

\subsection*{دو مدل ساده که نباید با هم اشتباه شوند}

نکته‌ای که اگر تصریح نشود شبیه تناقض به نظر می‌رسد: مدل ساده‌ی این پژوهش \lr{rMAE}
برابر \lr{0.849} دارد، یعنی ظاهراً مرجعی را که خودش تعریف می‌کند شکست می‌دهد.

توضیح آن است که دو مدل ساده‌ی متفاوت در کار است و هر دو استاندارد. مدل ساده‌ی گزارش‌شده
در جدول‌های فصل چهارم، قاعده‌ی «روز مشابه» است که \cref{sec:3-6} آن را تعریف می‌کند. اما مخرج
معیار \lr{rMAE}، بنا بر تعریف مرجع، پیش‌بینیِ ساده‌ی \emph{وقفه‌ی هفتگی} است: همان ساعت
در هفته‌ی پیش. پس \lr{0.849} یعنی قاعده‌ی روز مشابه حدود پانزده درصد از وقفه‌ی هفتگیِ
محض بهتر عمل می‌کند \lr{—} که دقیقاً همان چیزی است که باید بکند، و سنجه‌ای از ارزش
منطق رفتاریِ روز هفته.

\subsection*{چرا \lr{MAPE} گزارش نمی‌شود، و چرا \lr{sMAPE} با آن ناسازگار نیست}

از گزارش \lr{MAPE} متعارف آگاهانه صرف‌نظر شده است. در بازار برق آلمان قیمت‌های منفی
واقعاً رخ می‌دهند \lr{—} ۲۹۷ ساعت در داده‌ی آموزش و ۲۴۱ ساعت در دوره‌ی آزمون \lr{—} و
\lr{MAPE} بر خودِ $p_i$ تقسیم می‌کند؛ مخرجی که هم علامت‌پذیر است و هم در صفر
تعریف‌ناپذیر. یک ساعتِ نزدیک به صفر کافی است تا این معیار بدون کران منفجر شود.

اما \lr{sMAPE} گزارش می‌شود، و باید توضیح داد چرا این دو رفتار متفاوتی دارند. مخرج
\lr{sMAPE} میانگینِ \emph{قدرمطلق} مقدار واقعی و پیش‌بینی است، پس تنها در حالتی
مشکل‌ساز می‌شود که هر دو هم‌زمان به صفر نزدیک باشند \lr{—} حالتی به‌مراتب نادرتر.

شاهد تجربی این تفاوت در خودِ داده‌ی مرجع موجود است: بر پیش‌بینی‌های یکسان، مقاله‌ی مرجع
مقادیر \lr{MAPE} در بازه‌ی تقریبی ۷۷ تا ۱۳۷ درصد را در برابر \lr{sMAPE} در بازه‌ی ۱۴ تا
۱۷ درصد گزارش می‌کند \lr{—} حدود ده برابر \lr{—} و خود آن مقاله نیز \lr{MAPE} را در
این بازار نامطمئن می‌داند \cite{lago_forecasting_2021}.

با این‌همه، جمله‌ی تعیین‌کننده این است: \emph{هیچ نتیجه‌ای در این پایان‌نامه بر
\lr{sMAPE} استوار نیست.} تمام آزمون‌های معناداری بر زیان قدرمطلق محاسبه می‌شوند و هر
ادعای سرشناسی بر \lr{MAE} یا \lr{rMAE} بنا شده است؛ \lr{sMAPE} تنها برای امکان
مقایسه با جدول منتشرشده گزارش می‌گردد.

پیاده‌سازی هر چهار معیار مستقیماً از کتابخانه‌ی مرجع فراخوانی می‌شود و بازنویسی نشده
است، تا اعداد فصل چهارم بدون هیچ ابهام پیاده‌سازی با اعداد منتشرشده هم‌مقیاس بمانند.

\subsection*{آزمون معناداری دیبولد\lr{–}ماریانو}

برتری عددی یک مدل بر مدل دیگر ممکن است محصول نوفه باشد. برای هر زوج مدل، آزمون
دیبولد\lr{–}ماریانو\footnote{\lr{Diebold–Mariano}} بر سری تفاضل زیانِ روزانه اجرا
می‌شود. اگر $d_t$ تفاضل زیان دو مدل در روز $t$ باشد، آماره‌ی آزمون چنین است:

\begin{equation}
\mathrm{DM} = \frac{\bar{d}}{\sqrt{\widehat{\mathrm{Var}}_{\mathrm{HAC}}(\bar{d})}}
\end{equation}

که در آن $\bar{d}$ میانگین تفاضل زیان و مخرج، خطای معیارِ آن با تصحیح
\lr{HAC} است. این تصحیح ضروری است: خطاهای پیش‌بینی در روزهای متوالی همبسته‌اند و
نادیده‌گرفتن این همبستگی، معناداری را به‌صورت کاذب بزرگ نشان می‌دهد. تمام
\lr{p}-مقدارهای این پایان‌نامه یک‌طرفه و تصحیح‌شده‌اند، و هیچ‌جا مقدار تصحیح‌نشده نقل
نمی‌شود.

\subsection*{دو خانواده‌ی آزمون: تأییدی و اکتشافی}

\cref{sec:4-5-1} ماتریسی هفت در هفت از مقایسه‌های زوجی گزارش می‌کند، که ۲۱ آزمون مستقل
می‌شود. در چنین شماری، در سطح معناداری پنج درصد به‌طور میانگین حدود یک برنده‌ی کاذب
انتظار می‌رود. اعلام‌نکردن این نکته، هر نتیجه‌ی نزدیک به مرز را بی‌اعتبار می‌کند.

از این رو، پیش از آنکه هیچ \lr{p}-مقداری نقل شود، دو خانواده‌ی آزمون به‌صراحت تفکیک
می‌شوند.

\textbf{خانواده‌ی تأییدی، شامل یک فرضیه‌ی از پیش تعیین‌شده:} مقایسه‌ی وزن‌دهیِ رژیم‌محور
با وزن‌دهیِ ایستا. این فرضیه در دفتر تصمیمات پروژه ثبت شده است، در تاریخی که هنوز
مدل ترکیبیِ رژیم‌محور ساخته نشده و هیچ مدل ترکیبی‌ای بر دوره‌ی آزمون سنجیده نشده بود.
پس این فرضیه نمی‌توانسته با نگاه‌کردن به نتایج انتخاب شده باشد. یک فرضیه‌ی واحدِ از
پیش تعیین‌شده در سطح پنج درصد و بدون تصحیح آزموده می‌شود؛ این رویه‌ی استاندارد است و
مدخلِ تاریخ‌دار همان چیزی است که آن را قابل وارسی می‌کند.

\textbf{خانواده‌ی اکتشافی، شامل کل ماتریس ۲۱تایی:} برای این خانواده، هم
\lr{p}-مقدار خام و هم \lr{p}-مقدار تصحیح‌شده با روش هولم\lr{–}بونفرونی گزارش می‌شود.
گزارش‌کردن تنها مقدار خام نادرست است و پنهان‌کردن آن گریزگرایانه به نظر می‌رسد؛ پس
هر دو ستون می‌آید.

پیامد عملیِ این تفکیک در \cref{sec:4-5-1} دیده می‌شود و به سود پژوهش نیست: یکی از
تفاوت‌هایی که در نگاه نخست معنادار می‌نماید، پس از تصحیح دوام نمی‌آورد و در متن نیز
به‌عنوان تفاوتِ معنادار گزارش نمی‌شود.

\subsection*{تکرارپذیری}

بذر تصادفیِ سراسری ۴۲ در تمام اجزای تصادفی \lr{—} نمونه‌گیر جست‌وجوی ابرپارامتر،
مدل‌های درختی و شبکه‌های عصبی \lr{—} تثبیت شده و در کد به‌گونه‌ای تحمیل می‌شود که
هیچ فایل پیکربندی یا ابرپارامترِ تنظیم‌شده‌ای نتواند آن را بازنویسی کند. هر تصمیم
غیربدیهیِ روش‌شناختی نیز با تاریخ در دفتر ثبت تصمیمات پروژه مستند شده است، تا مسیر
رسیدن به هر انتخاب پس از پایان کار نیز قابل بازسازی باشد.
